\documentclass[11pt,a4paper]{article}
\usepackage[utf8]{inputenc}
\usepackage[T1]{fontenc}
\usepackage[english]{babel}
\usepackage{lmodern}
\usepackage{geometry}
\usepackage{xcolor}
\usepackage{tocloft}
\usepackage{hyperref}
\usepackage{enumitem}
\usepackage{parskip}
\usepackage{booktabs}
\usepackage{array}
\usepackage{longtable}
\usepackage{fancyhdr}
\usepackage{lastpage}
\usepackage{listings}
\usepackage{titlesec}
\geometry{margin=2.15cm}
\setlength{\headheight}{14pt}
\setlength{\emergencystretch}{3em}
\setlist{itemsep=3pt,parsep=1pt,topsep=4pt}
\definecolor{titleblue}{HTML}{123B5D}
\definecolor{softblue}{HTML}{EDF5FB}
\definecolor{softgray}{HTML}{F5F7FA}
\definecolor{accent}{HTML}{0F6E8C}
\hypersetup{
  colorlinks=true,linkcolor=titleblue,urlcolor=accent,
  pdftitle={Scientific Data Analysis v1.7: English publication edition},
  pdfauthor={interemi},
  pdfsubject={Real-world uses and scientific limits}
}
\titleformat{\section}{\Large\bfseries\color{titleblue}}{\thesection.}{0.6em}{}
\titleformat{\subsection}{\large\bfseries\color{accent}}{\thesubsection.}{0.6em}{}
\setlength{\cftsecnumwidth}{2.8em}
\tocloftpagestyle{fancy}
\lstdefinestyle{promptstyle}{
  basicstyle=\ttfamily\small,backgroundcolor=\color{softgray},
  frame=single,rulecolor=\color{softblue},breaklines=true,columns=fullflexible
}
\pagestyle{fancy}
\fancyhf{}
\fancyhead[L]{Scientific Data Analysis}
\fancyhead[R]{v1.7 guide: English edition}
\fancyfoot[C]{\thepage\ of \pageref{LastPage}}
\newcommand{\code}[1]{\texttt{\detokenize{#1}}}

\begin{document}
\pagenumbering{gobble}
\begin{titlepage}
\centering
\vspace*{1.0cm}
{\Huge\bfseries\color{titleblue} User Guide\\[0.15cm]
Scientific Data Analysis v1.7\par}
\vspace{0.65cm}
{\Large Real-world uses and scientific limits\par}
\vspace{0.8cm}
{\large Academic, professional, and personal workflows; 59 capabilities,
transferable patterns, and explicit applicability criteria\par}
\vspace{0.8cm}
{\large English publication edition: September 23, 2026\par}
\vspace{0.5cm}
\begin{minipage}{0.93\textwidth}
\small\raggedright
Translation of the preserved private v1.7 guide at commit \code{994eaea}.
Original TeX SHA-256:
\nolinkurl{a863704d5bfa264afee727339a83e752fab569e57e7e1f2c5cfb85af43c303e8}.
Private source:
\nolinkurl{skills/scientific-data-maintainer/.cache/active_references/guides/guia_scientific_data_analysis_v1_7.tex}.

Installation and validation claims describe the historical v1.7 record,
not the currently installed family. This translation does not rerun those
gates. Historical command examples use English illustrative input names;
script names, options, and contract identifiers are unchanged. Use fresh
output directories and safe input copies when reproducing examples.
The original remains unchanged privately.
\end{minipage}
\vfill
{\large Prepared with Codex\par}
{\large Copyright 2026 interemi\par}
\end{titlepage}

\clearpage
\pagenumbering{roman}
\begin{abstract}
\code{scientific-data-analysis} began with a strong focus on astronomy,
FITS, notebooks, LaTeX, data reduction, and scientific reporting. v1.7
extends its usefulness to academic, professional, and personal work while
preserving that foundation. A capability may be general-purpose, transfer
a pattern, remain domain-specific, be legacy, or serve maintainers only.
Analogies are not forced when a tool does not apply.

The original guide records the installed v1.7 state after final
synchronization and post-sync validation. The audited technical baseline
remains v1.6; v1.7 adds documentation, regressions, and focused fixes to
clarify what to use, what to avoid, and how to choose the next step.
\end{abstract}
\tableofcontents
\newpage
\pagenumbering{arabic}

\section{Executive overview}
v1.7 clarifies use beyond purely academic cases. It classifies all 59 public
capabilities, supplies general routes for tables, documents, notebooks,
containers, and time series, and explains transferable scientific patterns
and explicit reasons to reject a route.

Some tools work unchanged outside astronomy, some transfer only a pattern,
and others remain astrophysics, legacy, or maintenance tools. Accurately
describing those boundaries is part of the improvement.

\section{Historical product state}
\begin{itemize}[leftmargin=1.5em]
\item \textbf{v1.6 stable}: installed, audited baseline with 59 capabilities.
\item \textbf{v1.7 stable}: real-world extension, installed and validated after synchronization.
\item \textbf{Completed v1.7 phases}: classification, real-world intake, measurement/noise/QA, transferable astronomy patterns, professional/personal packages, documentation, release candidate, and final installation.
\item \textbf{State recorded by the source}: v1.7 installed; v1.6 remains the stable audit and regression baseline.
\end{itemize}

\section{Capability classification}
Maintained reference:
\begin{lstlisting}[style=promptstyle]
references/real-world-capability-classification-v1-7.md
\end{lstlisting}
\begin{longtable}{p{0.34\linewidth}r p{0.43\linewidth}}
\toprule
\textbf{Classification} & \textbf{Total} & \textbf{Meaning} \\
\midrule
\code{general} & 25 & Works outside astronomy without changing its conceptual contract. \\
\code{cross_domain_adaptable} & 9 & Originates in science, astronomy, or coursework, but transfers a real pattern with documentation, tests, and limits. \\
\code{domain_specific} & 7 & Useful in real settings within its original scientific domain. \\
\code{legacy_specific} & 12 & Depends on narrow tools, backends, practices, or platforms. \\
\code{maintainer_only} & 6 & Release, registry, smoke, probes, or internal validation. \\
\textbf{Total} & \textbf{59} & Matches the v1.6 public surface. \\
\bottomrule
\end{longtable}

\subsection{Reading the classification}
\begin{itemize}[leftmargin=1.5em]
\item \code{general}: a normal entry point for academic, business, or personal use.
\item \code{cross_domain_adaptable}: can inspire or execute a useful pattern; explain its original domain.
\item \code{domain_specific}: use only when input belongs to that domain.
\item \code{legacy_specific}: check environment and backend first; block cleanly when unsuitable.
\item \code{maintainer_only}: do not offer as an ordinary analysis tool.
\end{itemize}

\newpage
\section{General real-world routes}
\subsection{Tables and data folders}
For CSV, TSV, Excel, JSONL, parquet, or mixed folders:
\begin{lstlisting}[style=promptstyle]
python scripts/profile_table.py table.csv --summary-json profile.json
python scripts/cross_domain_data_workbench.py data_folder --output-dir out --summary-json out/summary.json
python scripts/duckdb_workbench.py table1.csv table2.csv --sql "SELECT * FROM source0 LIMIT 20"
\end{lstlisting}
Typical uses:
\begin{itemize}[leftmargin=1.5em]
\item Academic: inspect a course dataset before cleaning it.
\item Professional: tickets, events, sales, operations, or QA.
\item Personal: expenses, energy, habits, inventory, or historical exports.
\end{itemize}

\subsection{Documents and handoffs}
\begin{lstlisting}[style=promptstyle]
python scripts/document_intake_workbench.py document_package --output-dir intake
python scripts/semantic_diff.py baseline.csv candidate.csv --output-md diff.md
python scripts/deliverable_factory.py scaffold output --kind status-report --format markdown
\end{lstlisting}
These routes help identify folder contents, review priorities, differences
between versions, and the package to deliver.

\subsection{Inherited notebooks}
\begin{lstlisting}[style=promptstyle]
python scripts/notebook_workbench.py inspect notebook.ipynb --summary-json inspect.json
python scripts/notebook_workbench.py execute-copy notebook.ipynb --output-dir executed_copy
\end{lstlisting}
The purpose is copy-based execution, preservation of originals, risk
detection, and provenance. It is not automatic repair of every notebook.

\subsection{Containers}
\begin{lstlisting}[style=promptstyle]
python scripts/inspect_data_container.py package.zip --summary-json container.json
python scripts/inspect_data_container.py data.sqlite --output-json sqlite_inventory.json
\end{lstlisting}
Inspect before extracting or opening unfamiliar files. This is the intended
inspection entry point for ZIP, SQLite, NPZ, HDF5, MAT, NetCDF, parquet,
feather, and other supported containers; inspection is not a guarantee
that arbitrary input is safe.

\subsection{Time series}
\begin{lstlisting}[style=promptstyle]
python scripts/timeseries_forecasting_workbench.py \
  --output-dir forecast \
  --data-path series.csv \
  --date-column date \
  --value-column value \
  --frequency MS \
  --test-horizon 12
\end{lstlisting}
Output is a reproducible scaffold with a holdout, benchmarks, and clean
exports to start an analysis, not an infallible prediction.

\section{Transferable patterns}
\subsection{Measurement, noise, and QA}
References:
\begin{lstlisting}[style=promptstyle]
references/real-world-measurement-patterns.md
references/photometry-detectors.md
\end{lstlisting}
Patterns:
\begin{itemize}[leftmargin=1.5em]
\item Noise/SNR budgets for an ROI or sensor.
\item Linear instrument calibration using standards.
\item Lightweight QA for NaN/Inf, suspicious signs, negative errors, duplicate timestamps, or non-monotonic times.
\item Inspection of measurements and their uncertainties before modeling.
\end{itemize}

\subsection{Patterns originating in astronomy}
Reference: \path{references/astro-transferable-patterns.md}.
\begin{longtable}{>{\raggedright\arraybackslash}p{0.32\linewidth}>{\raggedright\arraybackslash}p{0.27\linewidth}>{\raggedright\arraybackslash}p{0.28\linewidth}}
\toprule
\textbf{Capability} & \textbf{Reusable pattern} & \textbf{Limit} \\
\midrule
\code{catalog crossmatch-sky} & Coordinate matching with angular tolerance. & Not general GIS; no CRS, polygons, or routes. \\
\code{fits_rgb_batch.py} & Multichannel image with a manifest and QA. & Remains oriented toward FITS and filters. \\
\code{rgb visual export} & Reproducible visual file with before/after evidence. & Display-only, not calibrated flux. \\
\code{astrometry preflight} & Preflight before an expensive or fragile backend. & The solver is astrometric, not general-purpose. \\
\code{RV validate-manifest} & Validate a manifest before reporting. & RV/Systemic schema, not a universal manifest. \\
\code{deck style-audit} & Visual, geometry, and editability QA on copied decks. & Does not replace final visual review. \\
\bottomrule
\end{longtable}

\section{Examples by setting}
\subsection{Academic}
\begin{itemize}[leftmargin=1.5em]
\item Review an instructor's notebook without changing it.
\item Inspect FITS, WCS, spectra, or scientific tables.
\item Create a LaTeX report prepared for Overleaf.
\item Compare branches of a group assignment.
\item Check whether a result includes units, uncertainty, and limits.
\end{itemize}
\subsection{Professional}
\begin{itemize}[leftmargin=1.5em]
\item Inventory operational CSV, JSONL, and Excel files.
\item Review a supplier's document package before a meeting.
\item Execute a copy of an inherited KPI notebook.
\item Inspect a technical ZIP containing SQLite and a README before extraction.
\item Build a forecasting scaffold for demand or workload.
\item Identify impossible sensor values before reporting metrics.
\end{itemize}
\subsection{Personal}
\begin{itemize}[leftmargin=1.5em]
\item Profile a household budget or exported expenses.
\item Organize anonymized administrative documents.
\item Execute a copy of a downloaded tutorial notebook.
\item Inspect a personal archive before extraction.
\item Create an energy-use or habits notebook.
\item Review household sensor data with unusual timestamps or impossible values.
\end{itemize}

\section{When a capability does not apply}
Use \code{NO_APLICA}, \code{BLOCKED_CONTROLADO}, or retain domain-specific
scope when:
\begin{itemize}[leftmargin=1.5em]
\item Input does not meet the contract's format, units, or assumptions.
\item An optional backend is absent and no accurate fallback exists.
\item Output would seem useful only by hiding its original domain.
\item A visual product would be mistaken for calibrated data.
\item A maintainer-only capability is proposed for an end-user task.
\item A legacy route requires an environment the user does not have.
\item An analogy outside the domain would offer marketing rather than practical help.
\end{itemize}

\subsection{Concrete examples}
\begin{itemize}[leftmargin=1.5em]
\item \code{li6708_equivalent_width_workbench.py} does not generalize beyond its narrow spectroscopy scope.
\item \code{astrometry_net_workbench.py verify-existing-wcs} does not validate non-celestial maps.
\item \code{stilts_workbench.py} is not a substitute for \code{duckdb_workbench.py}.
\item \code{rgb_visual_fits_export.py} creates traceable visual containers, not calibrated data.
\item \code{capability_probe_matrix.py} is a maintenance tool, not a normal user action.
\end{itemize}

\section{Historical installed-v1.7 validation}
Maintained regressions recorded in the original guide:
\begin{lstlisting}[style=promptstyle]
python scripts/audit_v1_7_real_world_classification_regression.py
python scripts/audit_v1_7_phase1_real_world_intake_regression.py
python scripts/audit_v1_7_phase2_measurement_noise_qa_regression.py
python scripts/audit_v1_7_phase3_astro_transfer_patterns_regression.py
python scripts/audit_v1_7_phase4_professional_personal_packages_regression.py
\end{lstlisting}
Documentation and smoke gates:
\begin{lstlisting}[style=promptstyle]
python scripts/sync_public_surface_docs.py --check
python scripts/portable_smoke_test.py --profile core
\end{lstlisting}

\section{Changes relative to v1.6}
\begin{itemize}[leftmargin=1.5em]
\item v1.6 hardened the public surface and completed the technical audit of 59 capabilities.
\item v1.7 explains who each capability serves and when to avoid it.
\item v1.7 expands professional and personal examples while retaining academic rigor.
\item Domain-specific and legacy boundaries remain; the skill does not become a generic business suite.
\item v1.7 adds real-package playbooks for tables, documents, notebooks, containers, time series, and sensors.
\end{itemize}

\section{Recorded release conclusion}
v1.7 clarifies how to choose a useful entry point, transfer patterns where
appropriate, preserve scientific contracts, and block cleanly when a route
does not apply. The skill remains scientific and technical while explaining
its practical role in coursework, professional projects, and personal tasks.
\end{document}
